\documentclass[10pt,a4paper]{amsart}
\usepackage[margin=19mm]{geometry}
\usepackage{amsmath,amssymb,mathtools}
\usepackage[scr=rsfs,cal=euler]{mathalfa}
\usepackage{xfrac}
\usepackage[unicode,hidelinks,bookmarks=false]{hyperref}
\newcommand{\ie}{i.e.\ }
\newcommand{\cf}{cf.\ }
\newcommand{\resp}{respectively\ }
\newcommand{\Subsection}{\subsection}
\providecommand{\nobreakdash}{\nobreak\hbox{-}}
\setlength{\emergencystretch}{2em}
\multlinegap0pt
\usepackage{bm}
\usepackage{bbm}
\usepackage{dsfont}
\usepackage{xspace}
\usepackage{cancel}
\usepackage{mathtools}
\usepackage{amsthm}
\makeatletter
\@ifundefined{prop}{\newtheorem{prop}{Prop}}{}
\makeatother
\DeclareMathOperator{\ar}{ar\!}
\DeclareMathOperator{\id}{id}
\newcommand{\cat}[1]{\mathcal{#1}}
\newcommand{\str}[1]{\mathbb{#1}}
\title[K\H{o}nig = Ramsey, A compactness lemma for Ramsey categorie]{K\H{o}nig = Ramsey, A compactness lemma for Ramsey categories}
\author{Maximilian Hadek}
\date{Source: arXiv:2508.11465v1 (CC BY 4.0)}
\begin{document}
\maketitle

\noindent\emph{Source and licence.} K\H{o}nig = Ramsey, A compactness lemma for Ramsey categories. Version arXiv:2508.11465v1 (CC BY 4.0), distributed under the Creative Commons Attribution 4.0 International licence, as recorded in the arXiv licence metadata for this exact version and shown on the abstract page https://arxiv.org/abs/2508.11465v1. Full rights notice: licenses/arxiv-2508.11465-CC-BY-4.0.txt.\par\smallskip

\noindent\textit{Editorial scope.} Counted unit: the Prop whose source label is \texttt{prop:definability} in the article (source0.tex lines 540--548, source label \texttt{prop:definability}), delivered together with the article's own complete original proof of it (source0.tex lines 550--577, one \texttt{proof} environment of 28 lines). No mathematics is written, repaired or re-derived: statement and proof are the article's own text line for line. Required context retained uncounted: none. Every definition, display and label of those lines is retained unchanged. Omitted from this package and not redistributed: 1--539, 549--549, 578--1100. Nothing inside a retained excerpt is omitted or altered: every retained body is delivered exactly as the article's lines are written, so no omission receipt of any kind accompanies this package. Citations: the retained excerpts carry no citation command at all, so no bibliography entry of the article is transcribed and no reference is redistributed. Reference adaptation: none was needed and none was made; every reference inside the delivered unit resolves inside it, so no unresolved reference or citation is produced. Printed slips kept literal: no printed word, symbol, hypothesis or number of the counted statement or of its proof was altered. Layout and class: the article's own class is \texttt{elsarticle} with packages including microtype, euscript, amsmath, amsfonts, amssymb, amsthm, tikz-cd, thmtools, hyperref; the wrapper is the lane's pinned \texttt{amsart} editorial wrapper at 10pt on A4, which restates the theorem-like environments the retained text needs and adds the article's own macro definitions that the retained text needs (\texttt{\textbackslash ar}, \texttt{\textbackslash cat}, \texttt{\textbackslash id}, \texttt{\textbackslash str}), with extra wrapper packages (bm, bbm, dsfont, xspace, cancel, mathtools). The delivered numbering is the unit's own. Assets: no retained span contains a figure environment, a graphics-inclusion command, a PDF-page inclusion, a table or a TikZ picture; no image file is copied into this package, the package declares no asset dependency and no image file is redistributed with it. Licence: version arXiv:2508.11465v1 of this article is distributed under the Creative Commons Attribution 4.0 International licence, as recorded in the arXiv licence metadata for this exact version and shown on the abstract page https://arxiv.org/abs/2508.11465v1. A full rights notice is delivered at licenses/arxiv-2508.11465-CC-BY-4.0.txt. Characters: every retained excerpt is pure ASCII, so the delivered engine, XeLaTeX with its default Latin Modern text font, reports no missing character.
\begin{prop}
\label{prop:definability}
    Let $\cat E\xrightarrow{\pi}\cat C$ be an abstract expansion between two classes of relational structures, where both $\cat C$ and $\cat E$ are closed under taking induced substructures, and $\pi$ preserves the underlying sets and maps. Let the signatures of $\cat C$ and $\cat E$ be $\sigma$ and $\tau$ respectively. If $\kappa$ is a regular cardinal larger than the number of symbols in $\sigma$ and the number of non-isomorphic structures in $\cat E$, then for every symbol $R\in\sigma$ there is a quantifier free $L_{\kappa,0}$-formula $\phi_R$ whose set of variables is the arity of $R$, which defines $\pi$, i.e.
    $$
    \pi(\str E) = (E; \phi_R^{\str E}\mid R\in\sigma)\in\cat C
    $$
    for all $\str E\in\cat E$.
    In particular, one can $L_{\kappa,0}$-define $\cat C$ from $\cat E$ without quantifiers.
\end{prop}

\begin{proof}
    For each structure $\str E\in\cat E$, let $\psi_{\cat E}$ be the formula on variable set $E$, which determines $\str E$, i.e.
    $$
    \psi_{\str E} = 
    \bigwedge_{\substack{S\in\tau\\ \bar{e}\in S^{\str E}}} S(\bar{e}) \wedge
    \bigwedge_{\substack{S\in\tau\\ \bar{e}\notin S^{\str E}}} \neg S(\bar{e})
    $$
    Given a $\tau$-structure $\str E'$ and an injective tuple $E\xrightarrow{f}E'$, the statement $\psi_{\str E}(f)$ will be true in $\str E'$ if and only if $f$ is an embedding from $\str E$ to $\str E'$. To deal with non-injective tuples, for every surjective map $N\xrightarrow{g}M$ we define the following formula with variable set $N$.
    $$
    \psi_g =
    \bigwedge_{\substack{x,y\in N\\ g(x)= g(y)}} (x = y) \wedge
    \bigwedge_{\substack{x,y\in N\\ g(x)\neq g(y)}} (x \neq y)
    $$
    For any symbol $R\in\sigma$, let $\phi_R$ be the following formula with variable set $\ar R$:
    $$
    \phi_R :=
    \bigvee_{\str E\in\cat E}
    \bigvee_{\substack{\bar{e}\in R^{\pi(\str E)} \\ \text{surj}}}
    \psi_{\bar{e}}\wedge \psi_{\str E}(\iota)
    $$
    Here, $\bar{e}$ is any surjective tuple the relation $R^{\pi(\str E)}$, while $\iota$ is any map $E\xrightarrow{\iota}\ar R$ which is left inverse to $\bar{e}$, so $\bar{e}\circ\iota=\id_{E}$. The size of the above formulae is bounded by the size of the signature $\tau$ and the number of non-isomorphic structures in $\cat E$ whose domain size is not larger than the arity of $S$.
    
    To verify that the formulae $\phi_R$ define $\pi$, let $\str E$ be a structure in $\cat E$ and let $\ar S\xrightarrow{\bar{c}} E$ be a tuple. We show that $\bar{c}\in R^{\pi(\str E)}$ if and only if $\phi_R(\bar{c})$ is true in $\str E$.
    
    Assume that $\phi_R(\bar{c})$ is true in $\str E$, then there is a structure $\str E'$ and a surjective tuple $\bar{e}\in R^{\pi(\str E')}$ such that $\psi_{\bar{e}}(\bar{c})$ and $\psi_{\str E}({\bar{c}\circ\iota})$ are true in $\str E'$. The former implies that $e_i=e_j$ if and only if $c_i=c_j$ and in particular that $\bar{c}\circ\iota$ is injective. Hence, it is an embedding $\str E'\xrightarrow{}\str E$ and also an embedding $\pi(\str E')\xrightarrow{}\pi(\str E)$. This implies that $\bar{c}=\bar{c}\circ\iota\circ\bar{e}$ is a tuple in $R^{\pi(\str E)}$.

    Conversely, if $\bar{c}$ is an element of $R^{\pi(\str E)}$, let $\str E'$ be the substructure of $\str E$ induced on the image of $\bar{c}$, then $\bar{c}$ is a surjective tuple in $R^{\pi(\str E')}$. Moreover, both $\psi_{\bar{c}}(\bar{c})$ and $\psi_{\str E'}(\bar{c}\circ\iota)$ hold in $\str E$, since $\bar{c}\circ\iota$ is the inclusion $\str E'\leq \str E$.
\end{proof}

\end{document}
